\documentclass[11pt,a4paper]{article}
\usepackage[T1]{fontenc}
\usepackage[utf8]{inputenc}
\usepackage[brazil]{babel}
\usepackage[top=1.5cm, bottom=1.5cm, left=2cm, right=2cm]{geometry}
\usepackage{hyperref}
\usepackage{enumitem}
\usepackage{titlesec}
\usepackage{xcolor}
\usepackage{parskip}
\usepackage{lmodern}
\renewcommand{\familydefault}{\sfdefault}

\definecolor{accent}{RGB}{59,44,73}

\pagestyle{empty}
\setlength{\parindent}{0pt}
\setlength{\parskip}{4pt}

\titleformat{\section}{\large\bfseries\color{accent}}{}{0em}{}[\color{accent}\titlerule]
\titlespacing{\section}{0pt}{10pt}{5pt}

\setlist[itemize]{noitemsep, topsep=2pt, leftmargin=1.2em}

\hypersetup{colorlinks=true, urlcolor=accent, hidelinks}

\begin{document}

% ── CABEÇALHO ──────────────────────────────────────────────────────────────────
{\LARGE\bfseries Samuel Louvera}\\[3pt]
{\large Estagiário de T.I \textbar{} Infraestrutura, Suporte e Desenvolvimento}\\[5pt]
Ibirité, MG \quad\textbullet\quad
\href{mailto:samuelouvera@gmail.com}{samuelouvera@gmail.com} \quad\textbullet\quad
(31) 99788-1386\\
\href{https://linkedin.com/in/samuel-louvera}{linkedin.com/in/samuel-louvera} \quad\textbullet\quad
\href{https://samuzaum.github.io/Portfolio/}{samuzaum.github.io/Portfolio/}

% ── RESUMO ─────────────────────────────────────────────────────────────────────
\section{Resumo}
Estudante de Análise e Desenvolvimento de Sistemas com formação técnica em T.I e atuação prática em
infraestrutura, suporte e desenvolvimento. Experiência com administração de identidades e acessos (IAM),
migração de dados para nuvem, suporte N1/N2 via GLPI/ITIL, manutenção de hardware e desenvolvimento
de soluções em Python, Laravel e Next.js para sistemas internos da empresa. Perfil proativo, organizado
e autônomo, com experiência em trabalho remoto e foco em resolução de problemas de forma independente.

% ── COMPETÊNCIAS ───────────────────────────────────────────────────────────────
\section{Competências Técnicas}
\begin{itemize}
  \item Administração de identidades e acessos (IAM) via Google Workspace
  \item Onboarding/offboarding em ambiente híbrido (nuvem e servidor local)
  \item Migração de dados e gestão de ciclo de vida de arquivos em Google Drive
  \item Backup, integridade de dados e políticas de recuperação de desastres
  \item Suporte técnico N1/N2 com GLPI e metodologia ITIL
  \item Manutenção de hardware: desktop, notebooks, firmwares e drivers
  \item Automação com Python: integração de dados e middleware
  \item Sistemas Operacionais: Windows, Linux (básico), macOS (básico)
  \item Configuração de redes: TCP/IP, roteadores, Wi-Fi
  \item Criação de manuais e documentação técnica para comunicação assíncrona em equipes remotas
  \item Ferramentas de colaboração remota: Google Meet, Microsoft Teams
  \item Virtualização: VirtualBox (criação e gerenciamento de máquinas virtuais)
  \item Ferramentas de IA: ChatGPT, Claude, Gemini, Grok (uso aplicado no dia a dia)
  \item Inglês técnico: leitura e interpretação de documentação
  \item Cisco Networking Academy: Network Defense e Endpoint Security
\end{itemize}

% ── EXPERIÊNCIA ────────────────────────────────────────────────────────────────
\section{Experiência}

\textbf{Lápis Raro} \hfill \textit{abr/2026 -- atual}\\
\textit{Estagiário de T.I}
\begin{itemize}
  \item Desenvolvimento e implementação de middleware em Python para automação de extração e
        integração de dados via protocolo Intelbras (em produção).
  \item Implementação de funcionalidades em sistemas internos da empresa (Laravel e Next.js/Supabase),
        incluindo um portal de gestão de documentos em produção e um dashboard financeiro em desenvolvimento.
  \item Administração de IAM via Google Workspace, cobrindo 95\% do parque de permissões com
        Princípio do Menor Privilégio.
  \item Gestão do ciclo de vida do colaborador: onboarding/offboarding em ambiente híbrido
        (nuvem e servidor local).
  \item Migração massiva de dados legados (CDs/HDs) para Google Drive com rotinas de backup
        e monitoramento de integridade.
  \item Suporte técnico N1/N2 via GLPI com metodologia ITIL: manutenção de hardware,
        atualização de firmware/drivers e refresh tecnológico.
\end{itemize}

\medskip

\textbf{AprendaNet \& Professor Global} \hfill \textit{ago/2025 -- atual · remoto}\\
\textit{Bolsista em Conteúdo Educacional (Projeto CNPq)}
\begin{itemize}
  \item Suporte técnico a usuários em ambiente LMS, de forma totalmente remota.
  \item Criação de manuais e procedimentos para colaboradores.
  \item Organização de informações e apoio à resolução de problemas em plataformas digitais.
\end{itemize}

\medskip

\textbf{Panpharma Distribuidora de Medicamentos} \hfill \textit{mar/2025 -- jun/2025}\\
\textit{Auxiliar de Logística}
\begin{itemize}
  \item Organização de processos operacionais e controle de informações.
  \item Apoio à equipe visando eficiência e redução de erros.
\end{itemize}

\medskip

\textbf{Polimac Equipamentos para Gastronomia} \hfill \textit{out/2024 -- jan/2025}\\
\textit{Estagiário de Vendas}
\begin{itemize}
  \item Atendimento ao cliente, suporte pós-venda e organização de pedidos.
  \item Gestão de e-commerce (Nuvem Shop e Mercado Livre).
\end{itemize}

% ── PROJETOS ───────────────────────────────────────────────────────────────────
\section{Projetos}

\textbf{Laboratório Prático de Infraestrutura} \textit{(experiência própria)}
\begin{itemize}
  \item Formatação e instalação de Windows em computadores pessoais.
  \item Montagem e manutenção de desktops: memória, HD/SSD e periféricos.
  \item Configuração de redes domésticas: roteador, Wi-Fi, IP.
\end{itemize}

\medskip

\textbf{Dashboard de Controle de Estoque}
\begin{itemize}
  \item Aplicação web com HTML, CSS, PHP e MySQL: CRUD, controle de acesso e
        gerenciamento de usuários.
\end{itemize}

% ── FORMAÇÃO ───────────────────────────────────────────────────────────────────
\section{Formação}

\textbf{Análise e Desenvolvimento de Sistemas}, Faculdade Descomplica \hfill \textit{2024 -- 2026}\\
\textbf{Técnico em Tecnologia da Informação}, SENAI BH CECOTEG \hfill \textit{2023 -- 2024}

% ── CERTIFICAÇÕES ──────────────────────────────────────────────────────────────
\section{Certificações}
\begin{itemize}
  \item Cisco Networking Academy: Network Defense
  \item Cisco Networking Academy: Endpoint Security
  \item Object-Oriented Developer
  \item Excel Básico
\end{itemize}

\end{document}
